\chapter{Limitations and Future Work}

\section{Limitations}

While the proposed RL-based routing framework demonstrates strong improvements in reliability and competitive performance, several limitations remain.

First, the evaluation is conducted using synthetic traffic patterns. Although these patterns provide controlled and diverse stress scenarios, they do not fully capture the complexity and temporal variability of real application workloads \cite{Malikov2025}. Second, the current formulation focuses primarily on reliability and latency. Other important system-level metrics, such as energy consumption and power efficiency \cite{Wu2017}, are not explicitly considered in the routing policy. Third, the study is limited to a 2D mesh NoC topology. Emerging architectures such as 2.5D and 3D NoCs \cite{Kondoth2026} introduce additional challenges, including increased thermal coupling and vertical interconnect constraints, which are not captured in this work. Finally, the proposed approach is evaluated in a simulation environment. While the use of a cycle-accurate simulator \cite{gem5} provides detailed insights, hardware-level validation is necessary to assess practical implementation overheads and feasibility.

\section{Future Work}

There are several promising directions for extending this work. A natural next step is to evaluate the proposed routing framework using real application workloads, such as PARSEC \cite{PARSEC} and SPLASH-2 \cite{Splash2}, to better understand its behavior under realistic traffic conditions.

The reinforcement learning formulation can be extended to incorporate additional objectives, such as energy efficiency, enabling a more comprehensive multi-objective optimization framework. In particular, replacing the fixed weighted-sum reward function with a dynamic or Pareto-based approach \cite{BahlousBoldi2025ParetoOptimal} would allow the routing policy to adaptively balance trade-offs between reliability and performance. Future work can also explore the application of this approach to emerging NoC architectures, including 2.5D and 3D designs \cite{Kondoth2026}, where reliability and thermal challenges are more pronounced. Finally, implementing and validating the proposed RL-based routing policy in hardware would provide valuable insights into its practical overheads, scalability, and real-world applicability. 

Together, these directions provide a pathway toward more adaptive, efficient, and reliable NoC routing strategies in future many-core systems.